% =====================================================================
%  eval_forms.tex  —  HUST Graduation-Thesis evaluation forms (Vietnamese)
%  Included from main.tex via % =====================================================================
%  eval_forms.tex  —  HUST Graduation-Thesis evaluation forms (Vietnamese)
%  Included from main.tex via % =====================================================================
%  eval_forms.tex  —  HUST Graduation-Thesis evaluation forms (Vietnamese)
%  Included from main.tex via % =====================================================================
%  eval_forms.tex  —  HUST Graduation-Thesis evaluation forms (Vietnamese)
%  Included from main.tex via \input{chapters/eval_forms}
%
%  This file is SELF-CONTAINED: it defines its own column types and a
%  local row-helper so it does not depend on macro definitions elsewhere.
%  It compiles with pdfLaTeX + [T5]{fontenc} (Vietnamese) and the packages
%  already loaded in main.tex (tabularx, booktabs/array, multirow,
%  xcolor[table], enumitem).  No siunitx / math is used here.
%
%  Student fields are read from the \student* macros defined in main.tex
%  (\studentnameVN, \studentid, \thesistitle).  Edit those in main.tex.
% =====================================================================

% ---- Local column types (multipliers sum to 4 = number of X columns) -
%  Column 1 (STT)     : narrow, centred
%  Column 2 (Tiêu chí): medium,  left
%  Column 3 (Hướng dẫn): wide,   left, justified
%  Column 4 (Điểm)    : narrow, centred
\newcolumntype{A}{>{\centering\arraybackslash\hsize=0.30\hsize}X}
\newcolumntype{B}{>{\raggedright\arraybackslash\hsize=0.95\hsize}X}
\newcolumntype{G}{>{\raggedright\arraybackslash\hsize=2.30\hsize}X}
\newcolumntype{P}{>{\centering\arraybackslash\hsize=0.45\hsize}X}

% ---- Local helpers ---------------------------------------------------
% Body font for the three forms. If any single form spills one line past
% its page on your machine, change \footnotesize -> \scriptsize here only.
\newcommand{\evalheaderfont}{\footnotesize}
% A grey header row for the four-column criteria tables:
\newcommand{\evalcolhead}{%
  \rowcolor{LightCyan}
  \textbf{STT} & \textbf{Tiêu chí (Điểm tối đa)} &
  \textbf{Hướng dẫn đánh giá tiêu chí} & \textbf{Điểm tiêu chí}\\
}

% =====================================================================
%  These admin forms are front matter: keep them single-spaced and
%  unnumbered, each on its own page, and excluded from the ToC.
% =====================================================================
\begingroup
\singlespacing
\setlength{\extrarowheight}{1pt}\setlength{\tabcolsep}{4pt}
\evalheaderfont

% =====================================================================
%  FORM 1 — DÀNH CHO CÁN BỘ HƯỚNG DẪN
% =====================================================================
\thispagestyle{fancy}
\begin{center}
  {\bfseries ĐẠI HỌC BÁCH KHOA HÀ NỘI}\\[2pt]
  {\bfseries TRƯỜNG ĐIỆN - ĐIỆN TỬ}\\[4pt]
  {\large\bfseries ĐÁNH GIÁ ĐỒ ÁN TỐT NGHIỆP}\\[2pt]
  {\bfseries (DÀNH CHO CÁN BỘ HƯỚNG DẪN)}
\end{center}

\vspace{3pt}
\noindent\textbf{Tên đề tài:} \thesistitle

\vspace{3pt}
\noindent\textbf{Họ tên SV:} \studentnameVN \hfill \textbf{MSSV:} \studentid

\vspace{2pt}
\noindent Cán bộ hướng dẫn 1: \dotfill

\vspace{2pt}
\noindent Cán bộ hướng dẫn 2: \dotfill

\vspace{3pt}
\noindent
\begin{tabularx}{\textwidth}{|A|B|G|P|}
  \hline
  \evalcolhead
  \hline
  \multirow{2}{=}{1}
    & \multirow{2}{=}{\textbf{Thái độ làm việc} (2,5 điểm)}
    & Nghiêm túc, tích cực và chủ động trong quá trình làm ĐATN & \\ \cline{3-4}
    & & Hoàn thành đầy đủ và đúng tiến độ các nội dung được GVHD giao & \\ \hline
  \multirow{2}{=}{2}
    & \multirow{2}{=}{\textbf{Kỹ năng viết quyển ĐATN} (2 điểm)}
    & Trình bày đúng mẫu quy định, bố cục các chương logic và hợp lý:
      bảng biểu, hình ảnh rõ ràng, có tiêu đề, được đánh số thứ tự và được
      giải thích hay đề cập đến trong đồ án, có căn lề, dấu cách sau dấu
      chấm, dấu phẩy, có mở đầu chương và kết luận chương, có liệt kê tài
      liệu tham khảo và có trích dẫn, v.v. & \\ \cline{3-4}
    & & Kỹ năng diễn đạt, phân tích, giải thích, lập luận: cấu trúc câu rõ
        ràng, văn phong khoa học, lập luận logic và có cơ sở, thuật ngữ
        chuyên ngành phù hợp, v.v. & \\ \hline
  \multirow{3}{=}{3}
    & \multirow{3}{=}{\textbf{Nội dung và kết quả đạt được} (5 điểm)}
    & Nêu rõ tính cấp thiết, ý nghĩa khoa học và thực tiễn của đề tài, các
      vấn đề và các giả thuyết, phạm vi ứng dụng của đề tài. Thực hiện đầy
      đủ quy trình nghiên cứu: đặt vấn đề, mục tiêu đề ra, phương pháp
      nghiên cứu/giải quyết vấn đề, kết quả đạt được, đánh giá và kết luận. & \\ \cline{3-4}
    & & Nội dung và kết quả được trình bày một cách logic và hợp lý, được
        phân tích và đánh giá thỏa đáng. Biện luận phân tích kết quả mô
        phỏng/phần mềm/thực nghiệm, so sánh kết quả đạt được với kết quả
        trước đó có liên quan. & \\ \cline{3-4}
    & & Chỉ rõ sự phù hợp giữa kết quả đạt được và mục tiêu ban đầu đề ra,
        đồng thời cung cấp lập luận để đề xuất hướng giải quyết có thể thực
        hiện trong tương lai. Hàm lượng khoa học/độ phức tạp cao, có tính
        mới/tính sáng tạo trong nội dung và kết quả đồ án. & \\ \hline
  \multirow{2}{=}{4}
    & \multirow{2}{=}{\textbf{Điểm thành tích} (1 điểm)}
    & Có bài báo KH được đăng hoặc chấp nhận đăng / đạt giải SV NCKH giải 3
      cấp Trường trở lên / các giải thưởng khoa học trong nước, quốc tế từ
      giải 3 trở lên / có đăng ký bằng phát minh sáng chế. (1 điểm) & \\ \cline{3-4}
    & & Được báo cáo tại hội đồng cấp Trường trong hội nghị SV NCKH nhưng
        không đạt giải từ giải 3 trở lên / đạt giải khuyến khích trong cuộc
        thi khoa học trong nước, quốc tế / kết quả đồ án là sản phẩm ứng
        dụng có tính hoàn thiện cao, yêu cầu khối lượng thực hiện lớn.
        (0,5 điểm) & \\ \hline
  \multicolumn{3}{|l|}{\textbf{Điểm tổng các tiêu chí:}} & \\ \hline
  \multicolumn{3}{|l|}{\textbf{Điểm hướng dẫn:}} & \\ \hline
\end{tabularx}

\vspace{3pt}
\hfill\begin{minipage}{6cm}\centering
  \textbf{Cán bộ hướng dẫn}\\
  (Ký và ghi rõ họ tên)
\end{minipage}

\vspace{2pt}\par\noindent\scriptsize\textit{Điểm từng tiêu chí cho lẻ đến 0,5. Nếu Điểm
tổng các tiêu chí $>$ 10 thì Điểm hướng dẫn làm tròn thành 10.}
\normalsize
\clearpage

% =====================================================================
%  FORM 2 — DÀNH CHO CÁN BỘ THÀNH VIÊN HỘI ĐỒNG
% =====================================================================
\thispagestyle{fancy}
\begin{center}
  {\bfseries ĐẠI HỌC BÁCH KHOA HÀ NỘI}\\[2pt]
  {\bfseries TRƯỜNG ĐIỆN - ĐIỆN TỬ}\\[4pt]
  {\large\bfseries ĐÁNH GIÁ ĐỒ ÁN TỐT NGHIỆP}\\[2pt]
  {\bfseries (DÀNH CHO CÁN BỘ THÀNH VIÊN HỘI ĐỒNG)}
\end{center}

\vspace{3pt}
\noindent Hội đồng số: \dotfill

\vspace{2pt}
\noindent\textbf{Họ tên SV:} \studentnameVN \hfill \textbf{MSSV:} \studentid

\vspace{2pt}
\noindent Cán bộ thành viên HĐ: \dotfill

\vspace{3pt}
\noindent
\begin{tabularx}{\textwidth}{|A|B|G|P|}
  \hline
  \evalcolhead
  \hline
  \multirow{2}{=}{1}
    & \multirow{2}{=}{\textbf{Chất lượng slides/Bản vẽ kỹ thuật} (1,5 điểm)}
    & Sử dụng các minh họa hỗ trợ: hình ảnh, biểu đồ rõ nét và phù hợp,
      dễ hiểu & \\ \cline{3-4}
    & & Không quá nhiều từ, biết sử dụng từ khoá; bố cục logic, có đánh số
        trang & \\ \hline
  \multirow{2}{=}{2}
    & \multirow{2}{=}{\textbf{Kỹ năng thuyết trình} (1,5 điểm)}
    & Tự tin, làm chủ nội dung trình bày, đúng thời gian quy định & \\ \cline{3-4}
    & & Dễ hiểu, dễ theo dõi, lô-gic, lôi cuốn. & \\ \hline
  \multirow{3}{=}{3}
    & \multirow{3}{=}{\textbf{Nội dung và kết quả đạt được} (5 điểm)}
    & Nêu rõ tính cấp thiết, ý nghĩa khoa học và thực tiễn của đề tài, các
      vấn đề và các giả thuyết, phạm vi ứng dụng của đề tài. Thực hiện đầy
      đủ quy trình nghiên cứu: đặt vấn đề, mục tiêu đề ra, phương pháp
      nghiên cứu/giải quyết vấn đề, kết quả đạt được, đánh giá và kết luận. & \\ \cline{3-4}
    & & Nội dung và kết quả được trình bày một cách logic và hợp lý, được
        phân tích và đánh giá thỏa đáng. Biện luận phân tích kết quả mô
        phỏng/phần mềm/thực nghiệm, so sánh kết quả đạt được với kết quả
        trước đó có liên quan. & \\ \cline{3-4}
    & & Chỉ rõ sự phù hợp giữa kết quả đạt được và mục tiêu ban đầu đề ra,
        đồng thời cung cấp lập luận để đề xuất hướng giải quyết có thể thực
        hiện trong tương lai. Hàm lượng khoa học/độ phức tạp cao, có tính
        mới/tính sáng tạo trong nội dung và kết quả đồ án. & \\ \hline
  \multirow{2}{=}{4}
    & \multirow{2}{=}{\textbf{Trả lời câu hỏi} (2,5 điểm)}
    & Trả lời ngắn gọn, chính xác, đi thẳng vào vấn đề của câu hỏi. & \\ \cline{3-4}
    & & Nắm vững kiến thức cơ bản liên quan đến lĩnh vực nghiên cứu/công
        việc của đồ án. & \\ \hline
  \multirow{2}{=}{5}
    & \multirow{2}{=}{\textbf{Điểm thành tích} (1 điểm)}
    & Có bài báo KH được đăng hoặc chấp nhận đăng / đạt giải SV NCKH giải 3
      cấp Trường trở lên / các giải thưởng khoa học trong nước, quốc tế từ
      giải 3 trở lên / có đăng ký bằng phát minh sáng chế. (1 điểm) & \\ \cline{3-4}
    & & Được báo cáo tại hội đồng cấp Trường trong hội nghị SV NCKH nhưng
        không đạt giải từ giải 3 trở lên / đạt giải khuyến khích trong cuộc
        thi khoa học trong nước, quốc tế / kết quả đồ án là sản phẩm ứng
        dụng có tính hoàn thiện cao, yêu cầu khối lượng thực hiện lớn.
        (0,5 điểm) & \\ \hline
  \multicolumn{3}{|l|}{\textbf{Điểm tổng các tiêu chí:}} & \\ \hline
  \multicolumn{3}{|l|}{\textbf{Điểm đánh giá:}} & \\ \hline
\end{tabularx}

\vspace{3pt}
\hfill\begin{minipage}{6cm}\centering
  \textbf{Cán bộ thành viên hội đồng}\\
  (Ký và ghi rõ họ tên)
\end{minipage}

\vspace{2pt}\par\noindent\scriptsize\textit{Điểm từng tiêu chí cho lẻ đến 0,5. Nếu Điểm
tổng các tiêu chí $>$ 10 thì Điểm đánh giá làm tròn thành 10.}
\normalsize
\clearpage

% =====================================================================
%  FORM 3 — DÀNH CHO CÁN BỘ PHẢN BIỆN
% =====================================================================
\thispagestyle{fancy}
\begin{center}
  {\bfseries ĐẠI HỌC BÁCH KHOA HÀ NỘI}\\[2pt]
  {\bfseries TRƯỜNG ĐIỆN - ĐIỆN TỬ}\\[4pt]
  {\large\bfseries ĐÁNH GIÁ ĐỒ ÁN TỐT NGHIỆP}\\[2pt]
  {\bfseries (DÀNH CHO CÁN BỘ PHẢN BIỆN)}
\end{center}

\vspace{3pt}
\noindent\textbf{Tên đề tài:} \thesistitle

\vspace{3pt}
\noindent\textbf{Họ tên SV:} \studentnameVN \hfill \textbf{MSSV:} \studentid

\vspace{2pt}
\noindent Cán bộ phản biện: \dotfill

\vspace{3pt}
\noindent
\begin{tabularx}{\textwidth}{|A|B|G|P|}
  \hline
  \evalcolhead
  \hline
  \multirow{2}{=}{1}
    & \multirow{2}{=}{\textbf{Trình bày quyển ĐATN} (4 điểm)}
    & Đồ án trình bày đúng mẫu quy định, bố cục các chương logic và hợp lý:
      bảng biểu, hình ảnh rõ ràng, có tiêu đề, được đánh số thứ tự và được
      giải thích hay đề cập đến trong đồ án, có căn lề, dấu cách sau dấu
      chấm, dấu phẩy, có mở đầu chương và kết luận chương, có liệt kê tài
      liệu tham khảo và có trích dẫn, v.v. & \\ \cline{3-4}
    & & Kỹ năng diễn đạt, phân tích, giải thích, lập luận: cấu trúc câu rõ
        ràng, văn phong khoa học, lập luận logic và có cơ sở, thuật ngữ
        chuyên ngành phù hợp, v.v. & \\ \hline
  \multirow{3}{=}{2}
    & \multirow{3}{=}{\textbf{Nội dung và kết quả đạt được} (5,5 điểm)}
    & Nêu rõ tính cấp thiết, ý nghĩa khoa học và thực tiễn của đề tài, các
      vấn đề và các giả thuyết, phạm vi ứng dụng của đề tài. Thực hiện đầy
      đủ quy trình nghiên cứu: đặt vấn đề, mục tiêu đề ra, phương pháp
      nghiên cứu/giải quyết vấn đề, kết quả đạt được, đánh giá và kết luận. & \\ \cline{3-4}
    & & Nội dung và kết quả được trình bày một cách logic và hợp lý, được
        phân tích và đánh giá thỏa đáng. Biện luận phân tích kết quả mô
        phỏng/phần mềm/thực nghiệm, so sánh kết quả đạt được với kết quả
        trước đó có liên quan. & \\ \cline{3-4}
    & & Chỉ rõ sự phù hợp giữa kết quả đạt được và mục tiêu ban đầu đề ra,
        đồng thời cung cấp lập luận để đề xuất hướng giải quyết có thể thực
        hiện trong tương lai. Hàm lượng khoa học/độ phức tạp cao, có tính
        mới/tính sáng tạo trong nội dung và kết quả đồ án. & \\ \hline
  \multirow{2}{=}{3}
    & \multirow{2}{=}{\textbf{Điểm thành tích} (1 điểm)}
    & Có bài báo KH được đăng hoặc chấp nhận đăng / đạt giải SV NCKH giải 3
      cấp Trường trở lên / các giải thưởng khoa học trong nước, quốc tế từ
      giải 3 trở lên / có đăng ký bằng phát minh sáng chế. (1 điểm) & \\ \cline{3-4}
    & & Được báo cáo tại hội đồng cấp Trường trong hội nghị SV NCKH nhưng
        không đạt giải từ giải 3 trở lên / đạt giải khuyến khích trong cuộc
        thi khoa học trong nước, quốc tế / kết quả đồ án là sản phẩm ứng
        dụng có tính hoàn thiện cao, yêu cầu khối lượng thực hiện lớn.
        (0,5 điểm) & \\ \hline
  \multicolumn{3}{|l|}{\textbf{Điểm tổng các tiêu chí:}} & \\ \hline
  \multicolumn{3}{|l|}{\textbf{Điểm phản biện:}} & \\ \hline
\end{tabularx}

\vspace{3pt}
\hfill\begin{minipage}{6cm}\centering
  \textbf{Cán bộ phản biện}\\
  (Ký và ghi rõ họ tên)
\end{minipage}

\vspace{2pt}\par\noindent\scriptsize\textit{Điểm từng tiêu chí cho lẻ đến 0,5. Nếu Điểm
tổng các tiêu chí $>$ 10 thì Điểm phản biện làm tròn thành 10.}
\normalsize

\endgroup
\clearpage
%
%  This file is SELF-CONTAINED: it defines its own column types and a
%  local row-helper so it does not depend on macro definitions elsewhere.
%  It compiles with pdfLaTeX + [T5]{fontenc} (Vietnamese) and the packages
%  already loaded in main.tex (tabularx, booktabs/array, multirow,
%  xcolor[table], enumitem).  No siunitx / math is used here.
%
%  Student fields are read from the \student* macros defined in main.tex
%  (\studentnameVN, \studentid, \thesistitle).  Edit those in main.tex.
% =====================================================================

% ---- Local column types (multipliers sum to 4 = number of X columns) -
%  Column 1 (STT)     : narrow, centred
%  Column 2 (Tiêu chí): medium,  left
%  Column 3 (Hướng dẫn): wide,   left, justified
%  Column 4 (Điểm)    : narrow, centred
\newcolumntype{A}{>{\centering\arraybackslash\hsize=0.30\hsize}X}
\newcolumntype{B}{>{\raggedright\arraybackslash\hsize=0.95\hsize}X}
\newcolumntype{G}{>{\raggedright\arraybackslash\hsize=2.30\hsize}X}
\newcolumntype{P}{>{\centering\arraybackslash\hsize=0.45\hsize}X}

% ---- Local helpers ---------------------------------------------------
% Body font for the three forms. If any single form spills one line past
% its page on your machine, change \footnotesize -> \scriptsize here only.
\newcommand{\evalheaderfont}{\footnotesize}
% A grey header row for the four-column criteria tables:
\newcommand{\evalcolhead}{%
  \rowcolor{LightCyan}
  \textbf{STT} & \textbf{Tiêu chí (Điểm tối đa)} &
  \textbf{Hướng dẫn đánh giá tiêu chí} & \textbf{Điểm tiêu chí}\\
}

% =====================================================================
%  These admin forms are front matter: keep them single-spaced and
%  unnumbered, each on its own page, and excluded from the ToC.
% =====================================================================
\begingroup
\singlespacing
\setlength{\extrarowheight}{1pt}\setlength{\tabcolsep}{4pt}
\evalheaderfont

% =====================================================================
%  FORM 1 — DÀNH CHO CÁN BỘ HƯỚNG DẪN
% =====================================================================
\thispagestyle{fancy}
\begin{center}
  {\bfseries ĐẠI HỌC BÁCH KHOA HÀ NỘI}\\[2pt]
  {\bfseries TRƯỜNG ĐIỆN - ĐIỆN TỬ}\\[4pt]
  {\large\bfseries ĐÁNH GIÁ ĐỒ ÁN TỐT NGHIỆP}\\[2pt]
  {\bfseries (DÀNH CHO CÁN BỘ HƯỚNG DẪN)}
\end{center}

\vspace{3pt}
\noindent\textbf{Tên đề tài:} \thesistitle

\vspace{3pt}
\noindent\textbf{Họ tên SV:} \studentnameVN \hfill \textbf{MSSV:} \studentid

\vspace{2pt}
\noindent Cán bộ hướng dẫn 1: \dotfill

\vspace{2pt}
\noindent Cán bộ hướng dẫn 2: \dotfill

\vspace{3pt}
\noindent
\begin{tabularx}{\textwidth}{|A|B|G|P|}
  \hline
  \evalcolhead
  \hline
  \multirow{2}{=}{1}
    & \multirow{2}{=}{\textbf{Thái độ làm việc} (2,5 điểm)}
    & Nghiêm túc, tích cực và chủ động trong quá trình làm ĐATN & \\ \cline{3-4}
    & & Hoàn thành đầy đủ và đúng tiến độ các nội dung được GVHD giao & \\ \hline
  \multirow{2}{=}{2}
    & \multirow{2}{=}{\textbf{Kỹ năng viết quyển ĐATN} (2 điểm)}
    & Trình bày đúng mẫu quy định, bố cục các chương logic và hợp lý:
      bảng biểu, hình ảnh rõ ràng, có tiêu đề, được đánh số thứ tự và được
      giải thích hay đề cập đến trong đồ án, có căn lề, dấu cách sau dấu
      chấm, dấu phẩy, có mở đầu chương và kết luận chương, có liệt kê tài
      liệu tham khảo và có trích dẫn, v.v. & \\ \cline{3-4}
    & & Kỹ năng diễn đạt, phân tích, giải thích, lập luận: cấu trúc câu rõ
        ràng, văn phong khoa học, lập luận logic và có cơ sở, thuật ngữ
        chuyên ngành phù hợp, v.v. & \\ \hline
  \multirow{3}{=}{3}
    & \multirow{3}{=}{\textbf{Nội dung và kết quả đạt được} (5 điểm)}
    & Nêu rõ tính cấp thiết, ý nghĩa khoa học và thực tiễn của đề tài, các
      vấn đề và các giả thuyết, phạm vi ứng dụng của đề tài. Thực hiện đầy
      đủ quy trình nghiên cứu: đặt vấn đề, mục tiêu đề ra, phương pháp
      nghiên cứu/giải quyết vấn đề, kết quả đạt được, đánh giá và kết luận. & \\ \cline{3-4}
    & & Nội dung và kết quả được trình bày một cách logic và hợp lý, được
        phân tích và đánh giá thỏa đáng. Biện luận phân tích kết quả mô
        phỏng/phần mềm/thực nghiệm, so sánh kết quả đạt được với kết quả
        trước đó có liên quan. & \\ \cline{3-4}
    & & Chỉ rõ sự phù hợp giữa kết quả đạt được và mục tiêu ban đầu đề ra,
        đồng thời cung cấp lập luận để đề xuất hướng giải quyết có thể thực
        hiện trong tương lai. Hàm lượng khoa học/độ phức tạp cao, có tính
        mới/tính sáng tạo trong nội dung và kết quả đồ án. & \\ \hline
  \multirow{2}{=}{4}
    & \multirow{2}{=}{\textbf{Điểm thành tích} (1 điểm)}
    & Có bài báo KH được đăng hoặc chấp nhận đăng / đạt giải SV NCKH giải 3
      cấp Trường trở lên / các giải thưởng khoa học trong nước, quốc tế từ
      giải 3 trở lên / có đăng ký bằng phát minh sáng chế. (1 điểm) & \\ \cline{3-4}
    & & Được báo cáo tại hội đồng cấp Trường trong hội nghị SV NCKH nhưng
        không đạt giải từ giải 3 trở lên / đạt giải khuyến khích trong cuộc
        thi khoa học trong nước, quốc tế / kết quả đồ án là sản phẩm ứng
        dụng có tính hoàn thiện cao, yêu cầu khối lượng thực hiện lớn.
        (0,5 điểm) & \\ \hline
  \multicolumn{3}{|l|}{\textbf{Điểm tổng các tiêu chí:}} & \\ \hline
  \multicolumn{3}{|l|}{\textbf{Điểm hướng dẫn:}} & \\ \hline
\end{tabularx}

\vspace{3pt}
\hfill\begin{minipage}{6cm}\centering
  \textbf{Cán bộ hướng dẫn}\\
  (Ký và ghi rõ họ tên)
\end{minipage}

\vspace{2pt}\par\noindent\scriptsize\textit{Điểm từng tiêu chí cho lẻ đến 0,5. Nếu Điểm
tổng các tiêu chí $>$ 10 thì Điểm hướng dẫn làm tròn thành 10.}
\normalsize
\clearpage

% =====================================================================
%  FORM 2 — DÀNH CHO CÁN BỘ THÀNH VIÊN HỘI ĐỒNG
% =====================================================================
\thispagestyle{fancy}
\begin{center}
  {\bfseries ĐẠI HỌC BÁCH KHOA HÀ NỘI}\\[2pt]
  {\bfseries TRƯỜNG ĐIỆN - ĐIỆN TỬ}\\[4pt]
  {\large\bfseries ĐÁNH GIÁ ĐỒ ÁN TỐT NGHIỆP}\\[2pt]
  {\bfseries (DÀNH CHO CÁN BỘ THÀNH VIÊN HỘI ĐỒNG)}
\end{center}

\vspace{3pt}
\noindent Hội đồng số: \dotfill

\vspace{2pt}
\noindent\textbf{Họ tên SV:} \studentnameVN \hfill \textbf{MSSV:} \studentid

\vspace{2pt}
\noindent Cán bộ thành viên HĐ: \dotfill

\vspace{3pt}
\noindent
\begin{tabularx}{\textwidth}{|A|B|G|P|}
  \hline
  \evalcolhead
  \hline
  \multirow{2}{=}{1}
    & \multirow{2}{=}{\textbf{Chất lượng slides/Bản vẽ kỹ thuật} (1,5 điểm)}
    & Sử dụng các minh họa hỗ trợ: hình ảnh, biểu đồ rõ nét và phù hợp,
      dễ hiểu & \\ \cline{3-4}
    & & Không quá nhiều từ, biết sử dụng từ khoá; bố cục logic, có đánh số
        trang & \\ \hline
  \multirow{2}{=}{2}
    & \multirow{2}{=}{\textbf{Kỹ năng thuyết trình} (1,5 điểm)}
    & Tự tin, làm chủ nội dung trình bày, đúng thời gian quy định & \\ \cline{3-4}
    & & Dễ hiểu, dễ theo dõi, lô-gic, lôi cuốn. & \\ \hline
  \multirow{3}{=}{3}
    & \multirow{3}{=}{\textbf{Nội dung và kết quả đạt được} (5 điểm)}
    & Nêu rõ tính cấp thiết, ý nghĩa khoa học và thực tiễn của đề tài, các
      vấn đề và các giả thuyết, phạm vi ứng dụng của đề tài. Thực hiện đầy
      đủ quy trình nghiên cứu: đặt vấn đề, mục tiêu đề ra, phương pháp
      nghiên cứu/giải quyết vấn đề, kết quả đạt được, đánh giá và kết luận. & \\ \cline{3-4}
    & & Nội dung và kết quả được trình bày một cách logic và hợp lý, được
        phân tích và đánh giá thỏa đáng. Biện luận phân tích kết quả mô
        phỏng/phần mềm/thực nghiệm, so sánh kết quả đạt được với kết quả
        trước đó có liên quan. & \\ \cline{3-4}
    & & Chỉ rõ sự phù hợp giữa kết quả đạt được và mục tiêu ban đầu đề ra,
        đồng thời cung cấp lập luận để đề xuất hướng giải quyết có thể thực
        hiện trong tương lai. Hàm lượng khoa học/độ phức tạp cao, có tính
        mới/tính sáng tạo trong nội dung và kết quả đồ án. & \\ \hline
  \multirow{2}{=}{4}
    & \multirow{2}{=}{\textbf{Trả lời câu hỏi} (2,5 điểm)}
    & Trả lời ngắn gọn, chính xác, đi thẳng vào vấn đề của câu hỏi. & \\ \cline{3-4}
    & & Nắm vững kiến thức cơ bản liên quan đến lĩnh vực nghiên cứu/công
        việc của đồ án. & \\ \hline
  \multirow{2}{=}{5}
    & \multirow{2}{=}{\textbf{Điểm thành tích} (1 điểm)}
    & Có bài báo KH được đăng hoặc chấp nhận đăng / đạt giải SV NCKH giải 3
      cấp Trường trở lên / các giải thưởng khoa học trong nước, quốc tế từ
      giải 3 trở lên / có đăng ký bằng phát minh sáng chế. (1 điểm) & \\ \cline{3-4}
    & & Được báo cáo tại hội đồng cấp Trường trong hội nghị SV NCKH nhưng
        không đạt giải từ giải 3 trở lên / đạt giải khuyến khích trong cuộc
        thi khoa học trong nước, quốc tế / kết quả đồ án là sản phẩm ứng
        dụng có tính hoàn thiện cao, yêu cầu khối lượng thực hiện lớn.
        (0,5 điểm) & \\ \hline
  \multicolumn{3}{|l|}{\textbf{Điểm tổng các tiêu chí:}} & \\ \hline
  \multicolumn{3}{|l|}{\textbf{Điểm đánh giá:}} & \\ \hline
\end{tabularx}

\vspace{3pt}
\hfill\begin{minipage}{6cm}\centering
  \textbf{Cán bộ thành viên hội đồng}\\
  (Ký và ghi rõ họ tên)
\end{minipage}

\vspace{2pt}\par\noindent\scriptsize\textit{Điểm từng tiêu chí cho lẻ đến 0,5. Nếu Điểm
tổng các tiêu chí $>$ 10 thì Điểm đánh giá làm tròn thành 10.}
\normalsize
\clearpage

% =====================================================================
%  FORM 3 — DÀNH CHO CÁN BỘ PHẢN BIỆN
% =====================================================================
\thispagestyle{fancy}
\begin{center}
  {\bfseries ĐẠI HỌC BÁCH KHOA HÀ NỘI}\\[2pt]
  {\bfseries TRƯỜNG ĐIỆN - ĐIỆN TỬ}\\[4pt]
  {\large\bfseries ĐÁNH GIÁ ĐỒ ÁN TỐT NGHIỆP}\\[2pt]
  {\bfseries (DÀNH CHO CÁN BỘ PHẢN BIỆN)}
\end{center}

\vspace{3pt}
\noindent\textbf{Tên đề tài:} \thesistitle

\vspace{3pt}
\noindent\textbf{Họ tên SV:} \studentnameVN \hfill \textbf{MSSV:} \studentid

\vspace{2pt}
\noindent Cán bộ phản biện: \dotfill

\vspace{3pt}
\noindent
\begin{tabularx}{\textwidth}{|A|B|G|P|}
  \hline
  \evalcolhead
  \hline
  \multirow{2}{=}{1}
    & \multirow{2}{=}{\textbf{Trình bày quyển ĐATN} (4 điểm)}
    & Đồ án trình bày đúng mẫu quy định, bố cục các chương logic và hợp lý:
      bảng biểu, hình ảnh rõ ràng, có tiêu đề, được đánh số thứ tự và được
      giải thích hay đề cập đến trong đồ án, có căn lề, dấu cách sau dấu
      chấm, dấu phẩy, có mở đầu chương và kết luận chương, có liệt kê tài
      liệu tham khảo và có trích dẫn, v.v. & \\ \cline{3-4}
    & & Kỹ năng diễn đạt, phân tích, giải thích, lập luận: cấu trúc câu rõ
        ràng, văn phong khoa học, lập luận logic và có cơ sở, thuật ngữ
        chuyên ngành phù hợp, v.v. & \\ \hline
  \multirow{3}{=}{2}
    & \multirow{3}{=}{\textbf{Nội dung và kết quả đạt được} (5,5 điểm)}
    & Nêu rõ tính cấp thiết, ý nghĩa khoa học và thực tiễn của đề tài, các
      vấn đề và các giả thuyết, phạm vi ứng dụng của đề tài. Thực hiện đầy
      đủ quy trình nghiên cứu: đặt vấn đề, mục tiêu đề ra, phương pháp
      nghiên cứu/giải quyết vấn đề, kết quả đạt được, đánh giá và kết luận. & \\ \cline{3-4}
    & & Nội dung và kết quả được trình bày một cách logic và hợp lý, được
        phân tích và đánh giá thỏa đáng. Biện luận phân tích kết quả mô
        phỏng/phần mềm/thực nghiệm, so sánh kết quả đạt được với kết quả
        trước đó có liên quan. & \\ \cline{3-4}
    & & Chỉ rõ sự phù hợp giữa kết quả đạt được và mục tiêu ban đầu đề ra,
        đồng thời cung cấp lập luận để đề xuất hướng giải quyết có thể thực
        hiện trong tương lai. Hàm lượng khoa học/độ phức tạp cao, có tính
        mới/tính sáng tạo trong nội dung và kết quả đồ án. & \\ \hline
  \multirow{2}{=}{3}
    & \multirow{2}{=}{\textbf{Điểm thành tích} (1 điểm)}
    & Có bài báo KH được đăng hoặc chấp nhận đăng / đạt giải SV NCKH giải 3
      cấp Trường trở lên / các giải thưởng khoa học trong nước, quốc tế từ
      giải 3 trở lên / có đăng ký bằng phát minh sáng chế. (1 điểm) & \\ \cline{3-4}
    & & Được báo cáo tại hội đồng cấp Trường trong hội nghị SV NCKH nhưng
        không đạt giải từ giải 3 trở lên / đạt giải khuyến khích trong cuộc
        thi khoa học trong nước, quốc tế / kết quả đồ án là sản phẩm ứng
        dụng có tính hoàn thiện cao, yêu cầu khối lượng thực hiện lớn.
        (0,5 điểm) & \\ \hline
  \multicolumn{3}{|l|}{\textbf{Điểm tổng các tiêu chí:}} & \\ \hline
  \multicolumn{3}{|l|}{\textbf{Điểm phản biện:}} & \\ \hline
\end{tabularx}

\vspace{3pt}
\hfill\begin{minipage}{6cm}\centering
  \textbf{Cán bộ phản biện}\\
  (Ký và ghi rõ họ tên)
\end{minipage}

\vspace{2pt}\par\noindent\scriptsize\textit{Điểm từng tiêu chí cho lẻ đến 0,5. Nếu Điểm
tổng các tiêu chí $>$ 10 thì Điểm phản biện làm tròn thành 10.}
\normalsize

\endgroup
\clearpage
%
%  This file is SELF-CONTAINED: it defines its own column types and a
%  local row-helper so it does not depend on macro definitions elsewhere.
%  It compiles with pdfLaTeX + [T5]{fontenc} (Vietnamese) and the packages
%  already loaded in main.tex (tabularx, booktabs/array, multirow,
%  xcolor[table], enumitem).  No siunitx / math is used here.
%
%  Student fields are read from the \student* macros defined in main.tex
%  (\studentnameVN, \studentid, \thesistitle).  Edit those in main.tex.
% =====================================================================

% ---- Local column types (multipliers sum to 4 = number of X columns) -
%  Column 1 (STT)     : narrow, centred
%  Column 2 (Tiêu chí): medium,  left
%  Column 3 (Hướng dẫn): wide,   left, justified
%  Column 4 (Điểm)    : narrow, centred
\newcolumntype{A}{>{\centering\arraybackslash\hsize=0.30\hsize}X}
\newcolumntype{B}{>{\raggedright\arraybackslash\hsize=0.95\hsize}X}
\newcolumntype{G}{>{\raggedright\arraybackslash\hsize=2.30\hsize}X}
\newcolumntype{P}{>{\centering\arraybackslash\hsize=0.45\hsize}X}

% ---- Local helpers ---------------------------------------------------
% Body font for the three forms. If any single form spills one line past
% its page on your machine, change \footnotesize -> \scriptsize here only.
\newcommand{\evalheaderfont}{\footnotesize}
% A grey header row for the four-column criteria tables:
\newcommand{\evalcolhead}{%
  \rowcolor{LightCyan}
  \textbf{STT} & \textbf{Tiêu chí (Điểm tối đa)} &
  \textbf{Hướng dẫn đánh giá tiêu chí} & \textbf{Điểm tiêu chí}\\
}

% =====================================================================
%  These admin forms are front matter: keep them single-spaced and
%  unnumbered, each on its own page, and excluded from the ToC.
% =====================================================================
\begingroup
\singlespacing
\setlength{\extrarowheight}{1pt}\setlength{\tabcolsep}{4pt}
\evalheaderfont

% =====================================================================
%  FORM 1 — DÀNH CHO CÁN BỘ HƯỚNG DẪN
% =====================================================================
\thispagestyle{fancy}
\begin{center}
  {\bfseries ĐẠI HỌC BÁCH KHOA HÀ NỘI}\\[2pt]
  {\bfseries TRƯỜNG ĐIỆN - ĐIỆN TỬ}\\[4pt]
  {\large\bfseries ĐÁNH GIÁ ĐỒ ÁN TỐT NGHIỆP}\\[2pt]
  {\bfseries (DÀNH CHO CÁN BỘ HƯỚNG DẪN)}
\end{center}

\vspace{3pt}
\noindent\textbf{Tên đề tài:} \thesistitle

\vspace{3pt}
\noindent\textbf{Họ tên SV:} \studentnameVN \hfill \textbf{MSSV:} \studentid

\vspace{2pt}
\noindent Cán bộ hướng dẫn 1: \dotfill

\vspace{2pt}
\noindent Cán bộ hướng dẫn 2: \dotfill

\vspace{3pt}
\noindent
\begin{tabularx}{\textwidth}{|A|B|G|P|}
  \hline
  \evalcolhead
  \hline
  \multirow{2}{=}{1}
    & \multirow{2}{=}{\textbf{Thái độ làm việc} (2,5 điểm)}
    & Nghiêm túc, tích cực và chủ động trong quá trình làm ĐATN & \\ \cline{3-4}
    & & Hoàn thành đầy đủ và đúng tiến độ các nội dung được GVHD giao & \\ \hline
  \multirow{2}{=}{2}
    & \multirow{2}{=}{\textbf{Kỹ năng viết quyển ĐATN} (2 điểm)}
    & Trình bày đúng mẫu quy định, bố cục các chương logic và hợp lý:
      bảng biểu, hình ảnh rõ ràng, có tiêu đề, được đánh số thứ tự và được
      giải thích hay đề cập đến trong đồ án, có căn lề, dấu cách sau dấu
      chấm, dấu phẩy, có mở đầu chương và kết luận chương, có liệt kê tài
      liệu tham khảo và có trích dẫn, v.v. & \\ \cline{3-4}
    & & Kỹ năng diễn đạt, phân tích, giải thích, lập luận: cấu trúc câu rõ
        ràng, văn phong khoa học, lập luận logic và có cơ sở, thuật ngữ
        chuyên ngành phù hợp, v.v. & \\ \hline
  \multirow{3}{=}{3}
    & \multirow{3}{=}{\textbf{Nội dung và kết quả đạt được} (5 điểm)}
    & Nêu rõ tính cấp thiết, ý nghĩa khoa học và thực tiễn của đề tài, các
      vấn đề và các giả thuyết, phạm vi ứng dụng của đề tài. Thực hiện đầy
      đủ quy trình nghiên cứu: đặt vấn đề, mục tiêu đề ra, phương pháp
      nghiên cứu/giải quyết vấn đề, kết quả đạt được, đánh giá và kết luận. & \\ \cline{3-4}
    & & Nội dung và kết quả được trình bày một cách logic và hợp lý, được
        phân tích và đánh giá thỏa đáng. Biện luận phân tích kết quả mô
        phỏng/phần mềm/thực nghiệm, so sánh kết quả đạt được với kết quả
        trước đó có liên quan. & \\ \cline{3-4}
    & & Chỉ rõ sự phù hợp giữa kết quả đạt được và mục tiêu ban đầu đề ra,
        đồng thời cung cấp lập luận để đề xuất hướng giải quyết có thể thực
        hiện trong tương lai. Hàm lượng khoa học/độ phức tạp cao, có tính
        mới/tính sáng tạo trong nội dung và kết quả đồ án. & \\ \hline
  \multirow{2}{=}{4}
    & \multirow{2}{=}{\textbf{Điểm thành tích} (1 điểm)}
    & Có bài báo KH được đăng hoặc chấp nhận đăng / đạt giải SV NCKH giải 3
      cấp Trường trở lên / các giải thưởng khoa học trong nước, quốc tế từ
      giải 3 trở lên / có đăng ký bằng phát minh sáng chế. (1 điểm) & \\ \cline{3-4}
    & & Được báo cáo tại hội đồng cấp Trường trong hội nghị SV NCKH nhưng
        không đạt giải từ giải 3 trở lên / đạt giải khuyến khích trong cuộc
        thi khoa học trong nước, quốc tế / kết quả đồ án là sản phẩm ứng
        dụng có tính hoàn thiện cao, yêu cầu khối lượng thực hiện lớn.
        (0,5 điểm) & \\ \hline
  \multicolumn{3}{|l|}{\textbf{Điểm tổng các tiêu chí:}} & \\ \hline
  \multicolumn{3}{|l|}{\textbf{Điểm hướng dẫn:}} & \\ \hline
\end{tabularx}

\vspace{3pt}
\hfill\begin{minipage}{6cm}\centering
  \textbf{Cán bộ hướng dẫn}\\
  (Ký và ghi rõ họ tên)
\end{minipage}

\vspace{2pt}\par\noindent\scriptsize\textit{Điểm từng tiêu chí cho lẻ đến 0,5. Nếu Điểm
tổng các tiêu chí $>$ 10 thì Điểm hướng dẫn làm tròn thành 10.}
\normalsize
\clearpage

% =====================================================================
%  FORM 2 — DÀNH CHO CÁN BỘ THÀNH VIÊN HỘI ĐỒNG
% =====================================================================
\thispagestyle{fancy}
\begin{center}
  {\bfseries ĐẠI HỌC BÁCH KHOA HÀ NỘI}\\[2pt]
  {\bfseries TRƯỜNG ĐIỆN - ĐIỆN TỬ}\\[4pt]
  {\large\bfseries ĐÁNH GIÁ ĐỒ ÁN TỐT NGHIỆP}\\[2pt]
  {\bfseries (DÀNH CHO CÁN BỘ THÀNH VIÊN HỘI ĐỒNG)}
\end{center}

\vspace{3pt}
\noindent Hội đồng số: \dotfill

\vspace{2pt}
\noindent\textbf{Họ tên SV:} \studentnameVN \hfill \textbf{MSSV:} \studentid

\vspace{2pt}
\noindent Cán bộ thành viên HĐ: \dotfill

\vspace{3pt}
\noindent
\begin{tabularx}{\textwidth}{|A|B|G|P|}
  \hline
  \evalcolhead
  \hline
  \multirow{2}{=}{1}
    & \multirow{2}{=}{\textbf{Chất lượng slides/Bản vẽ kỹ thuật} (1,5 điểm)}
    & Sử dụng các minh họa hỗ trợ: hình ảnh, biểu đồ rõ nét và phù hợp,
      dễ hiểu & \\ \cline{3-4}
    & & Không quá nhiều từ, biết sử dụng từ khoá; bố cục logic, có đánh số
        trang & \\ \hline
  \multirow{2}{=}{2}
    & \multirow{2}{=}{\textbf{Kỹ năng thuyết trình} (1,5 điểm)}
    & Tự tin, làm chủ nội dung trình bày, đúng thời gian quy định & \\ \cline{3-4}
    & & Dễ hiểu, dễ theo dõi, lô-gic, lôi cuốn. & \\ \hline
  \multirow{3}{=}{3}
    & \multirow{3}{=}{\textbf{Nội dung và kết quả đạt được} (5 điểm)}
    & Nêu rõ tính cấp thiết, ý nghĩa khoa học và thực tiễn của đề tài, các
      vấn đề và các giả thuyết, phạm vi ứng dụng của đề tài. Thực hiện đầy
      đủ quy trình nghiên cứu: đặt vấn đề, mục tiêu đề ra, phương pháp
      nghiên cứu/giải quyết vấn đề, kết quả đạt được, đánh giá và kết luận. & \\ \cline{3-4}
    & & Nội dung và kết quả được trình bày một cách logic và hợp lý, được
        phân tích và đánh giá thỏa đáng. Biện luận phân tích kết quả mô
        phỏng/phần mềm/thực nghiệm, so sánh kết quả đạt được với kết quả
        trước đó có liên quan. & \\ \cline{3-4}
    & & Chỉ rõ sự phù hợp giữa kết quả đạt được và mục tiêu ban đầu đề ra,
        đồng thời cung cấp lập luận để đề xuất hướng giải quyết có thể thực
        hiện trong tương lai. Hàm lượng khoa học/độ phức tạp cao, có tính
        mới/tính sáng tạo trong nội dung và kết quả đồ án. & \\ \hline
  \multirow{2}{=}{4}
    & \multirow{2}{=}{\textbf{Trả lời câu hỏi} (2,5 điểm)}
    & Trả lời ngắn gọn, chính xác, đi thẳng vào vấn đề của câu hỏi. & \\ \cline{3-4}
    & & Nắm vững kiến thức cơ bản liên quan đến lĩnh vực nghiên cứu/công
        việc của đồ án. & \\ \hline
  \multirow{2}{=}{5}
    & \multirow{2}{=}{\textbf{Điểm thành tích} (1 điểm)}
    & Có bài báo KH được đăng hoặc chấp nhận đăng / đạt giải SV NCKH giải 3
      cấp Trường trở lên / các giải thưởng khoa học trong nước, quốc tế từ
      giải 3 trở lên / có đăng ký bằng phát minh sáng chế. (1 điểm) & \\ \cline{3-4}
    & & Được báo cáo tại hội đồng cấp Trường trong hội nghị SV NCKH nhưng
        không đạt giải từ giải 3 trở lên / đạt giải khuyến khích trong cuộc
        thi khoa học trong nước, quốc tế / kết quả đồ án là sản phẩm ứng
        dụng có tính hoàn thiện cao, yêu cầu khối lượng thực hiện lớn.
        (0,5 điểm) & \\ \hline
  \multicolumn{3}{|l|}{\textbf{Điểm tổng các tiêu chí:}} & \\ \hline
  \multicolumn{3}{|l|}{\textbf{Điểm đánh giá:}} & \\ \hline
\end{tabularx}

\vspace{3pt}
\hfill\begin{minipage}{6cm}\centering
  \textbf{Cán bộ thành viên hội đồng}\\
  (Ký và ghi rõ họ tên)
\end{minipage}

\vspace{2pt}\par\noindent\scriptsize\textit{Điểm từng tiêu chí cho lẻ đến 0,5. Nếu Điểm
tổng các tiêu chí $>$ 10 thì Điểm đánh giá làm tròn thành 10.}
\normalsize
\clearpage

% =====================================================================
%  FORM 3 — DÀNH CHO CÁN BỘ PHẢN BIỆN
% =====================================================================
\thispagestyle{fancy}
\begin{center}
  {\bfseries ĐẠI HỌC BÁCH KHOA HÀ NỘI}\\[2pt]
  {\bfseries TRƯỜNG ĐIỆN - ĐIỆN TỬ}\\[4pt]
  {\large\bfseries ĐÁNH GIÁ ĐỒ ÁN TỐT NGHIỆP}\\[2pt]
  {\bfseries (DÀNH CHO CÁN BỘ PHẢN BIỆN)}
\end{center}

\vspace{3pt}
\noindent\textbf{Tên đề tài:} \thesistitle

\vspace{3pt}
\noindent\textbf{Họ tên SV:} \studentnameVN \hfill \textbf{MSSV:} \studentid

\vspace{2pt}
\noindent Cán bộ phản biện: \dotfill

\vspace{3pt}
\noindent
\begin{tabularx}{\textwidth}{|A|B|G|P|}
  \hline
  \evalcolhead
  \hline
  \multirow{2}{=}{1}
    & \multirow{2}{=}{\textbf{Trình bày quyển ĐATN} (4 điểm)}
    & Đồ án trình bày đúng mẫu quy định, bố cục các chương logic và hợp lý:
      bảng biểu, hình ảnh rõ ràng, có tiêu đề, được đánh số thứ tự và được
      giải thích hay đề cập đến trong đồ án, có căn lề, dấu cách sau dấu
      chấm, dấu phẩy, có mở đầu chương và kết luận chương, có liệt kê tài
      liệu tham khảo và có trích dẫn, v.v. & \\ \cline{3-4}
    & & Kỹ năng diễn đạt, phân tích, giải thích, lập luận: cấu trúc câu rõ
        ràng, văn phong khoa học, lập luận logic và có cơ sở, thuật ngữ
        chuyên ngành phù hợp, v.v. & \\ \hline
  \multirow{3}{=}{2}
    & \multirow{3}{=}{\textbf{Nội dung và kết quả đạt được} (5,5 điểm)}
    & Nêu rõ tính cấp thiết, ý nghĩa khoa học và thực tiễn của đề tài, các
      vấn đề và các giả thuyết, phạm vi ứng dụng của đề tài. Thực hiện đầy
      đủ quy trình nghiên cứu: đặt vấn đề, mục tiêu đề ra, phương pháp
      nghiên cứu/giải quyết vấn đề, kết quả đạt được, đánh giá và kết luận. & \\ \cline{3-4}
    & & Nội dung và kết quả được trình bày một cách logic và hợp lý, được
        phân tích và đánh giá thỏa đáng. Biện luận phân tích kết quả mô
        phỏng/phần mềm/thực nghiệm, so sánh kết quả đạt được với kết quả
        trước đó có liên quan. & \\ \cline{3-4}
    & & Chỉ rõ sự phù hợp giữa kết quả đạt được và mục tiêu ban đầu đề ra,
        đồng thời cung cấp lập luận để đề xuất hướng giải quyết có thể thực
        hiện trong tương lai. Hàm lượng khoa học/độ phức tạp cao, có tính
        mới/tính sáng tạo trong nội dung và kết quả đồ án. & \\ \hline
  \multirow{2}{=}{3}
    & \multirow{2}{=}{\textbf{Điểm thành tích} (1 điểm)}
    & Có bài báo KH được đăng hoặc chấp nhận đăng / đạt giải SV NCKH giải 3
      cấp Trường trở lên / các giải thưởng khoa học trong nước, quốc tế từ
      giải 3 trở lên / có đăng ký bằng phát minh sáng chế. (1 điểm) & \\ \cline{3-4}
    & & Được báo cáo tại hội đồng cấp Trường trong hội nghị SV NCKH nhưng
        không đạt giải từ giải 3 trở lên / đạt giải khuyến khích trong cuộc
        thi khoa học trong nước, quốc tế / kết quả đồ án là sản phẩm ứng
        dụng có tính hoàn thiện cao, yêu cầu khối lượng thực hiện lớn.
        (0,5 điểm) & \\ \hline
  \multicolumn{3}{|l|}{\textbf{Điểm tổng các tiêu chí:}} & \\ \hline
  \multicolumn{3}{|l|}{\textbf{Điểm phản biện:}} & \\ \hline
\end{tabularx}

\vspace{3pt}
\hfill\begin{minipage}{6cm}\centering
  \textbf{Cán bộ phản biện}\\
  (Ký và ghi rõ họ tên)
\end{minipage}

\vspace{2pt}\par\noindent\scriptsize\textit{Điểm từng tiêu chí cho lẻ đến 0,5. Nếu Điểm
tổng các tiêu chí $>$ 10 thì Điểm phản biện làm tròn thành 10.}
\normalsize

\endgroup
\clearpage
%
%  This file is SELF-CONTAINED: it defines its own column types and a
%  local row-helper so it does not depend on macro definitions elsewhere.
%  It compiles with pdfLaTeX + [T5]{fontenc} (Vietnamese) and the packages
%  already loaded in main.tex (tabularx, booktabs/array, multirow,
%  xcolor[table], enumitem).  No siunitx / math is used here.
%
%  Student fields are read from the \student* macros defined in main.tex
%  (\studentnameVN, \studentid, \thesistitle).  Edit those in main.tex.
% =====================================================================

% ---- Local column types (multipliers sum to 4 = number of X columns) -
%  Column 1 (STT)     : narrow, centred
%  Column 2 (Tiêu chí): medium,  left
%  Column 3 (Hướng dẫn): wide,   left, justified
%  Column 4 (Điểm)    : narrow, centred
\newcolumntype{A}{>{\centering\arraybackslash\hsize=0.30\hsize}X}
\newcolumntype{B}{>{\raggedright\arraybackslash\hsize=0.95\hsize}X}
\newcolumntype{G}{>{\raggedright\arraybackslash\hsize=2.30\hsize}X}
\newcolumntype{P}{>{\centering\arraybackslash\hsize=0.45\hsize}X}

% ---- Local helpers ---------------------------------------------------
% Body font for the three forms. If any single form spills one line past
% its page on your machine, change \footnotesize -> \scriptsize here only.
\newcommand{\evalheaderfont}{\footnotesize}
% A grey header row for the four-column criteria tables:
\newcommand{\evalcolhead}{%
  \rowcolor{LightCyan}
  \textbf{STT} & \textbf{Tiêu chí (Điểm tối đa)} &
  \textbf{Hướng dẫn đánh giá tiêu chí} & \textbf{Điểm tiêu chí}\\
}

% =====================================================================
%  These admin forms are front matter: keep them single-spaced and
%  unnumbered, each on its own page, and excluded from the ToC.
% =====================================================================
\begingroup
\singlespacing
\setlength{\extrarowheight}{1pt}\setlength{\tabcolsep}{4pt}
\evalheaderfont

% =====================================================================
%  FORM 1 — DÀNH CHO CÁN BỘ HƯỚNG DẪN
% =====================================================================
\thispagestyle{fancy}
\begin{center}
  {\bfseries ĐẠI HỌC BÁCH KHOA HÀ NỘI}\\[2pt]
  {\bfseries TRƯỜNG ĐIỆN - ĐIỆN TỬ}\\[4pt]
  {\large\bfseries ĐÁNH GIÁ ĐỒ ÁN TỐT NGHIỆP}\\[2pt]
  {\bfseries (DÀNH CHO CÁN BỘ HƯỚNG DẪN)}
\end{center}

\vspace{3pt}
\noindent\textbf{Tên đề tài:} \thesistitle

\vspace{3pt}
\noindent\textbf{Họ tên SV:} \studentnameVN \hfill \textbf{MSSV:} \studentid

\vspace{2pt}
\noindent Cán bộ hướng dẫn 1: \dotfill

\vspace{2pt}
\noindent Cán bộ hướng dẫn 2: \dotfill

\vspace{3pt}
\noindent
\begin{tabularx}{\textwidth}{|A|B|G|P|}
  \hline
  \evalcolhead
  \hline
  \multirow{2}{=}{1}
    & \multirow{2}{=}{\textbf{Thái độ làm việc} (2,5 điểm)}
    & Nghiêm túc, tích cực và chủ động trong quá trình làm ĐATN & \\ \cline{3-4}
    & & Hoàn thành đầy đủ và đúng tiến độ các nội dung được GVHD giao & \\ \hline
  \multirow{2}{=}{2}
    & \multirow{2}{=}{\textbf{Kỹ năng viết quyển ĐATN} (2 điểm)}
    & Trình bày đúng mẫu quy định, bố cục các chương logic và hợp lý:
      bảng biểu, hình ảnh rõ ràng, có tiêu đề, được đánh số thứ tự và được
      giải thích hay đề cập đến trong đồ án, có căn lề, dấu cách sau dấu
      chấm, dấu phẩy, có mở đầu chương và kết luận chương, có liệt kê tài
      liệu tham khảo và có trích dẫn, v.v. & \\ \cline{3-4}
    & & Kỹ năng diễn đạt, phân tích, giải thích, lập luận: cấu trúc câu rõ
        ràng, văn phong khoa học, lập luận logic và có cơ sở, thuật ngữ
        chuyên ngành phù hợp, v.v. & \\ \hline
  \multirow{3}{=}{3}
    & \multirow{3}{=}{\textbf{Nội dung và kết quả đạt được} (5 điểm)}
    & Nêu rõ tính cấp thiết, ý nghĩa khoa học và thực tiễn của đề tài, các
      vấn đề và các giả thuyết, phạm vi ứng dụng của đề tài. Thực hiện đầy
      đủ quy trình nghiên cứu: đặt vấn đề, mục tiêu đề ra, phương pháp
      nghiên cứu/giải quyết vấn đề, kết quả đạt được, đánh giá và kết luận. & \\ \cline{3-4}
    & & Nội dung và kết quả được trình bày một cách logic và hợp lý, được
        phân tích và đánh giá thỏa đáng. Biện luận phân tích kết quả mô
        phỏng/phần mềm/thực nghiệm, so sánh kết quả đạt được với kết quả
        trước đó có liên quan. & \\ \cline{3-4}
    & & Chỉ rõ sự phù hợp giữa kết quả đạt được và mục tiêu ban đầu đề ra,
        đồng thời cung cấp lập luận để đề xuất hướng giải quyết có thể thực
        hiện trong tương lai. Hàm lượng khoa học/độ phức tạp cao, có tính
        mới/tính sáng tạo trong nội dung và kết quả đồ án. & \\ \hline
  \multirow{2}{=}{4}
    & \multirow{2}{=}{\textbf{Điểm thành tích} (1 điểm)}
    & Có bài báo KH được đăng hoặc chấp nhận đăng / đạt giải SV NCKH giải 3
      cấp Trường trở lên / các giải thưởng khoa học trong nước, quốc tế từ
      giải 3 trở lên / có đăng ký bằng phát minh sáng chế. (1 điểm) & \\ \cline{3-4}
    & & Được báo cáo tại hội đồng cấp Trường trong hội nghị SV NCKH nhưng
        không đạt giải từ giải 3 trở lên / đạt giải khuyến khích trong cuộc
        thi khoa học trong nước, quốc tế / kết quả đồ án là sản phẩm ứng
        dụng có tính hoàn thiện cao, yêu cầu khối lượng thực hiện lớn.
        (0,5 điểm) & \\ \hline
  \multicolumn{3}{|l|}{\textbf{Điểm tổng các tiêu chí:}} & \\ \hline
  \multicolumn{3}{|l|}{\textbf{Điểm hướng dẫn:}} & \\ \hline
\end{tabularx}

\vspace{3pt}
\hfill\begin{minipage}{6cm}\centering
  \textbf{Cán bộ hướng dẫn}\\
  (Ký và ghi rõ họ tên)
\end{minipage}

\vspace{2pt}\par\noindent\scriptsize\textit{Điểm từng tiêu chí cho lẻ đến 0,5. Nếu Điểm
tổng các tiêu chí $>$ 10 thì Điểm hướng dẫn làm tròn thành 10.}
\normalsize
\clearpage

% =====================================================================
%  FORM 2 — DÀNH CHO CÁN BỘ THÀNH VIÊN HỘI ĐỒNG
% =====================================================================
\thispagestyle{fancy}
\begin{center}
  {\bfseries ĐẠI HỌC BÁCH KHOA HÀ NỘI}\\[2pt]
  {\bfseries TRƯỜNG ĐIỆN - ĐIỆN TỬ}\\[4pt]
  {\large\bfseries ĐÁNH GIÁ ĐỒ ÁN TỐT NGHIỆP}\\[2pt]
  {\bfseries (DÀNH CHO CÁN BỘ THÀNH VIÊN HỘI ĐỒNG)}
\end{center}

\vspace{3pt}
\noindent Hội đồng số: \dotfill

\vspace{2pt}
\noindent\textbf{Họ tên SV:} \studentnameVN \hfill \textbf{MSSV:} \studentid

\vspace{2pt}
\noindent Cán bộ thành viên HĐ: \dotfill

\vspace{3pt}
\noindent
\begin{tabularx}{\textwidth}{|A|B|G|P|}
  \hline
  \evalcolhead
  \hline
  \multirow{2}{=}{1}
    & \multirow{2}{=}{\textbf{Chất lượng slides/Bản vẽ kỹ thuật} (1,5 điểm)}
    & Sử dụng các minh họa hỗ trợ: hình ảnh, biểu đồ rõ nét và phù hợp,
      dễ hiểu & \\ \cline{3-4}
    & & Không quá nhiều từ, biết sử dụng từ khoá; bố cục logic, có đánh số
        trang & \\ \hline
  \multirow{2}{=}{2}
    & \multirow{2}{=}{\textbf{Kỹ năng thuyết trình} (1,5 điểm)}
    & Tự tin, làm chủ nội dung trình bày, đúng thời gian quy định & \\ \cline{3-4}
    & & Dễ hiểu, dễ theo dõi, lô-gic, lôi cuốn. & \\ \hline
  \multirow{3}{=}{3}
    & \multirow{3}{=}{\textbf{Nội dung và kết quả đạt được} (5 điểm)}
    & Nêu rõ tính cấp thiết, ý nghĩa khoa học và thực tiễn của đề tài, các
      vấn đề và các giả thuyết, phạm vi ứng dụng của đề tài. Thực hiện đầy
      đủ quy trình nghiên cứu: đặt vấn đề, mục tiêu đề ra, phương pháp
      nghiên cứu/giải quyết vấn đề, kết quả đạt được, đánh giá và kết luận. & \\ \cline{3-4}
    & & Nội dung và kết quả được trình bày một cách logic và hợp lý, được
        phân tích và đánh giá thỏa đáng. Biện luận phân tích kết quả mô
        phỏng/phần mềm/thực nghiệm, so sánh kết quả đạt được với kết quả
        trước đó có liên quan. & \\ \cline{3-4}
    & & Chỉ rõ sự phù hợp giữa kết quả đạt được và mục tiêu ban đầu đề ra,
        đồng thời cung cấp lập luận để đề xuất hướng giải quyết có thể thực
        hiện trong tương lai. Hàm lượng khoa học/độ phức tạp cao, có tính
        mới/tính sáng tạo trong nội dung và kết quả đồ án. & \\ \hline
  \multirow{2}{=}{4}
    & \multirow{2}{=}{\textbf{Trả lời câu hỏi} (2,5 điểm)}
    & Trả lời ngắn gọn, chính xác, đi thẳng vào vấn đề của câu hỏi. & \\ \cline{3-4}
    & & Nắm vững kiến thức cơ bản liên quan đến lĩnh vực nghiên cứu/công
        việc của đồ án. & \\ \hline
  \multirow{2}{=}{5}
    & \multirow{2}{=}{\textbf{Điểm thành tích} (1 điểm)}
    & Có bài báo KH được đăng hoặc chấp nhận đăng / đạt giải SV NCKH giải 3
      cấp Trường trở lên / các giải thưởng khoa học trong nước, quốc tế từ
      giải 3 trở lên / có đăng ký bằng phát minh sáng chế. (1 điểm) & \\ \cline{3-4}
    & & Được báo cáo tại hội đồng cấp Trường trong hội nghị SV NCKH nhưng
        không đạt giải từ giải 3 trở lên / đạt giải khuyến khích trong cuộc
        thi khoa học trong nước, quốc tế / kết quả đồ án là sản phẩm ứng
        dụng có tính hoàn thiện cao, yêu cầu khối lượng thực hiện lớn.
        (0,5 điểm) & \\ \hline
  \multicolumn{3}{|l|}{\textbf{Điểm tổng các tiêu chí:}} & \\ \hline
  \multicolumn{3}{|l|}{\textbf{Điểm đánh giá:}} & \\ \hline
\end{tabularx}

\vspace{3pt}
\hfill\begin{minipage}{6cm}\centering
  \textbf{Cán bộ thành viên hội đồng}\\
  (Ký và ghi rõ họ tên)
\end{minipage}

\vspace{2pt}\par\noindent\scriptsize\textit{Điểm từng tiêu chí cho lẻ đến 0,5. Nếu Điểm
tổng các tiêu chí $>$ 10 thì Điểm đánh giá làm tròn thành 10.}
\normalsize
\clearpage

% =====================================================================
%  FORM 3 — DÀNH CHO CÁN BỘ PHẢN BIỆN
% =====================================================================
\thispagestyle{fancy}
\begin{center}
  {\bfseries ĐẠI HỌC BÁCH KHOA HÀ NỘI}\\[2pt]
  {\bfseries TRƯỜNG ĐIỆN - ĐIỆN TỬ}\\[4pt]
  {\large\bfseries ĐÁNH GIÁ ĐỒ ÁN TỐT NGHIỆP}\\[2pt]
  {\bfseries (DÀNH CHO CÁN BỘ PHẢN BIỆN)}
\end{center}

\vspace{3pt}
\noindent\textbf{Tên đề tài:} \thesistitle

\vspace{3pt}
\noindent\textbf{Họ tên SV:} \studentnameVN \hfill \textbf{MSSV:} \studentid

\vspace{2pt}
\noindent Cán bộ phản biện: \dotfill

\vspace{3pt}
\noindent
\begin{tabularx}{\textwidth}{|A|B|G|P|}
  \hline
  \evalcolhead
  \hline
  \multirow{2}{=}{1}
    & \multirow{2}{=}{\textbf{Trình bày quyển ĐATN} (4 điểm)}
    & Đồ án trình bày đúng mẫu quy định, bố cục các chương logic và hợp lý:
      bảng biểu, hình ảnh rõ ràng, có tiêu đề, được đánh số thứ tự và được
      giải thích hay đề cập đến trong đồ án, có căn lề, dấu cách sau dấu
      chấm, dấu phẩy, có mở đầu chương và kết luận chương, có liệt kê tài
      liệu tham khảo và có trích dẫn, v.v. & \\ \cline{3-4}
    & & Kỹ năng diễn đạt, phân tích, giải thích, lập luận: cấu trúc câu rõ
        ràng, văn phong khoa học, lập luận logic và có cơ sở, thuật ngữ
        chuyên ngành phù hợp, v.v. & \\ \hline
  \multirow{3}{=}{2}
    & \multirow{3}{=}{\textbf{Nội dung và kết quả đạt được} (5,5 điểm)}
    & Nêu rõ tính cấp thiết, ý nghĩa khoa học và thực tiễn của đề tài, các
      vấn đề và các giả thuyết, phạm vi ứng dụng của đề tài. Thực hiện đầy
      đủ quy trình nghiên cứu: đặt vấn đề, mục tiêu đề ra, phương pháp
      nghiên cứu/giải quyết vấn đề, kết quả đạt được, đánh giá và kết luận. & \\ \cline{3-4}
    & & Nội dung và kết quả được trình bày một cách logic và hợp lý, được
        phân tích và đánh giá thỏa đáng. Biện luận phân tích kết quả mô
        phỏng/phần mềm/thực nghiệm, so sánh kết quả đạt được với kết quả
        trước đó có liên quan. & \\ \cline{3-4}
    & & Chỉ rõ sự phù hợp giữa kết quả đạt được và mục tiêu ban đầu đề ra,
        đồng thời cung cấp lập luận để đề xuất hướng giải quyết có thể thực
        hiện trong tương lai. Hàm lượng khoa học/độ phức tạp cao, có tính
        mới/tính sáng tạo trong nội dung và kết quả đồ án. & \\ \hline
  \multirow{2}{=}{3}
    & \multirow{2}{=}{\textbf{Điểm thành tích} (1 điểm)}
    & Có bài báo KH được đăng hoặc chấp nhận đăng / đạt giải SV NCKH giải 3
      cấp Trường trở lên / các giải thưởng khoa học trong nước, quốc tế từ
      giải 3 trở lên / có đăng ký bằng phát minh sáng chế. (1 điểm) & \\ \cline{3-4}
    & & Được báo cáo tại hội đồng cấp Trường trong hội nghị SV NCKH nhưng
        không đạt giải từ giải 3 trở lên / đạt giải khuyến khích trong cuộc
        thi khoa học trong nước, quốc tế / kết quả đồ án là sản phẩm ứng
        dụng có tính hoàn thiện cao, yêu cầu khối lượng thực hiện lớn.
        (0,5 điểm) & \\ \hline
  \multicolumn{3}{|l|}{\textbf{Điểm tổng các tiêu chí:}} & \\ \hline
  \multicolumn{3}{|l|}{\textbf{Điểm phản biện:}} & \\ \hline
\end{tabularx}

\vspace{3pt}
\hfill\begin{minipage}{6cm}\centering
  \textbf{Cán bộ phản biện}\\
  (Ký và ghi rõ họ tên)
\end{minipage}

\vspace{2pt}\par\noindent\scriptsize\textit{Điểm từng tiêu chí cho lẻ đến 0,5. Nếu Điểm
tổng các tiêu chí $>$ 10 thì Điểm phản biện làm tròn thành 10.}
\normalsize

\endgroup
\clearpage